\chapter{Professional Ethics, Licensure, and Course Policies}

\section{Professional Engineering Licensure}
CLO~6 requires you to define the process for attaining professional engineering licensure and explain why it is important. The SDI Professional Development modules on licensure and liability (Week~2) cover this in detail; this section provides the reference framework.

\subsection{The Path to PE}
\begin{enumerate}[nosep,leftmargin=1.5em]
\item \textbf{Accredited degree}: earn an ABET-accredited engineering degree~[25]. You are completing this requirement now.
\item \textbf{FE exam}: pass the Fundamentals of Engineering exam, administered by NCEES~[20]. The FE is typically taken during your senior year or shortly after graduation. It is a 5-hour-and-20-minute, 110-question computer-based exam covering mathematics, sciences, and engineering fundamentals. Passing the FE makes you an \textbf{Engineer Intern (EI)} or \textbf{Engineer-in-Training (EIT)}.
\item \textbf{Supervised experience}: accumulate 4 years (in most states) of progressive engineering experience under the supervision of a licensed PE. The experience must be verified by your supervisors.
\item \textbf{PE exam}: pass the Principles and Practice of Engineering exam in your specific discipline (Mechanical, Civil, Electrical, Environmental, etc.). Passing makes you a \textbf{Professional Engineer (PE)} licensed in that state.
\end{enumerate}

\subsection{Why Licensure Matters}
A PE license is \textbf{legally required} to: offer engineering services directly to the public, sign and seal engineering documents (drawings, plans, specifications), and hold certain positions in government and consulting. Beyond legal requirements, licensure signals competence, ethical commitment, and personal accountability to the profession and the public. Many employers in consulting, construction, and government require or prefer PE licensure. The licensure process ensures that engineers offering services to the public have demonstrated a minimum level of competency and ethical commitment.

\subsection{Involuntary vs. Voluntary Liability}
The SDI Week~2 lecture introduces these legal concepts:

\textbf{Involuntary liability} (tort law): you are liable for harm caused by your engineering work even without a contract. If a product you designed injures someone, the injured party can sue for negligence. The standard is: would a reasonably competent engineer have foreseen and prevented this hazard?

\textbf{Voluntary liability} (contract law): you accept specific obligations when you sign a contract to provide engineering services. Failure to meet contractual obligations exposes you to breach-of-contract claims.

Understanding liability is not about being scared; it is about understanding why professional standards, thorough documentation, and safety margins exist. Every engineering decision you make in capstone (from material selection to safety factor to test procedure) reflects the same judgment that licensed engineers exercise daily.

\section{NCEES Rules of Professional Conduct}
The NCEES Model Rules~[20] define the ethical obligations of licensed engineers. These rules apply to you as aspiring professionals, even before licensure. The key principles are summarized below; this is not the complete text of the Model Rules. For the full, authoritative document, consult the NCEES \textit{Model Rules of Professional Conduct} at \url{https://ncees.org/licensure/laws-and-rules/}.

\subsection{Section A: Obligations to the Public}
Your first and highest obligation is to the \textbf{health, safety, and welfare of the public}~[19],~[20]. This is not a platitude; it is a legally enforceable professional duty. If your professional judgment is overruled and public safety is endangered, you have an obligation to report the situation to the appropriate authorities. You shall include all relevant information in professional reports and statements and shall not knowingly permit the publication of misleading information.

\subsection{Section B: Obligations to Employers and Clients}
Undertake only assignments for which you are \textbf{qualified by education or experience}. If a project requires skills you do not have, inform your client or employer. Do not reveal confidential information without the client's consent (NDA compliance; see Chapter~17). Disclose all \textbf{conflicts of interest}; situations where your personal interests could bias your professional judgment.

\subsection{Section C: Obligations to Other Licensees}
Do not \textbf{misrepresent your qualifications}. Do not injure the reputation of other licensees. Report known violations of the rules of professional conduct and material errors that may impact public safety.

\section{Engineering Ethics in Practice}
Ethics in engineering is not abstract philosophy. It is the daily practice of making decisions that affect real people, often under time pressure, incomplete information, and competing demands from stakeholders with different priorities.

\subsection{The Ethics Assignment}
The Ethics Assignment (SDI, Sprint~4) is a \textbf{1000--1500 word argumentative essay}. You must take a controversial stance on an ethical issue directly connected to your project's hardware, software, or system, defend it with evidence, and address opposing viewpoints. See Chapter~5 for essay structure and Chapter~18 for writing guidance.

The assignment requires you to move beyond generic ethical principles (``safety is important'') to specific, nuanced dilemmas found in your actual project. Because AI tools can easily generate generic ethical statements (e.g., ``Safety is important''), a top-tier essay must demonstrate \textbf{specific, nuanced critical thinking} that connects directly to the unique constraints of your Senior Design project.

\subsection{Case Studies}
\textbf{Challenger (1986)}: Morton Thiokol engineers warned that O-ring seals would fail at low temperatures. Management overruled the recommendation. Seven crew members died (Presidential Commission on the Space Shuttle Challenger Accident, 1986). Lesson: the obligation to public safety overrides schedule pressure and management hierarchy.

\textbf{Flint Water Crisis (2014--)}: cost-cutting decisions to change water sources without adequate corrosion control exposed a predominantly low-income, majority-Black community to lead contamination. Engineers who raised concerns were overruled or ignored (Michigan Civil Rights Commission, 2017). Lesson: equity and environmental justice are engineering responsibilities, not just political considerations.

\textbf{Boeing 737 MAX (2018--2019)}: design decisions to minimize pilot training requirements for a new flight control system (MCAS) contributed to two crashes and 346 deaths (Joint Authorities Technical Review, 2019). Lesson: when safety-critical systems are designed to minimize cost rather than maximize safety, the consequences can be catastrophic.

\textbf{Hyatt Regency Walkway Collapse (1981)}: a structural connection detail was changed during construction without adequate engineering review. 114 people died (NBS, 1982). Lesson: every change to a safety-critical design must be reviewed by a qualified engineer, no matter how minor it appears.

These case studies share a common thread: in each case, the technical failure was preceded by a \emph{process} failure; suppressed dissent, inadequate review, conflicted incentives, or organizational pressure to prioritize cost and schedule over safety. The ethical engineer's job is to recognize and resist these process failures.

\section{Generative AI Policy}
Students must represent all work submitted for evaluation as their own intellectual product. The course AI policy:

\textbf{Permitted uses}: studying, brainstorming, proofreading, exploring AI capabilities, generating initial research directions, checking grammar and clarity.

\textbf{Prohibited uses}: submitting AI-generated text, code, images, or solutions as your own work without explicit instructor permission and proper attribution. ``Hallucinated'' AI citations (references that do not actually exist) constitute an \textbf{academic integrity violation} and will be treated as such.

The purpose of this policy is not to prevent you from using AI tools; it is to ensure that the work you submit reflects \emph{your} engineering judgment, \emph{your} analysis, and \emph{your} understanding. An AI can generate a requirements document; but only \emph{you} can verify that those requirements accurately capture your client's needs and your project's constraints.

\section{Attendance Policy}
Attendance at lectures, recitation, and all team meetings is mandatory. Capstone Design is a synchronous course by design; your PA, client, and teammates depend on your presence.

\begin{center}\small
\begin{tabularx}{\textwidth}{cX}
\toprule
\textbf{Unexcused Absences} & \textbf{Consequence} \\
\midrule
1--3 & No grade deduction, but grade of 0 for any missed in-class work \\
4 & Final average lowered by 2 percentage points; formal PIP may be initiated \\
5--6 & Final grade lowered by 5 percentage points \\
7 or more & Final grade of 0 assigned; withdrawal is recommended \\
\bottomrule
\end{tabularx}
\end{center}

The Scrum Master is encouraged to track attendance at all team meetings. Absences must be documented with dates, times, and context to support any PIP action (Chapter~13). It is not enough to say ``student X missed several meetings''; you need dates and the record from the Communication Lead's attendance log.

\section{Late Work and Grading}
Late individual assignments: \textbf{10\% penalty per 24 hours} late, at PA discretion. No individual assignments accepted after 5 calendar days. Team assignments: discuss extensions with your PA before the deadline; retroactive extensions are not granted.

Grades returned by PAs are preliminary; the CF may apply correction factors for consistency across sections. Grade scale: A (93.34--100\%), A$-$ (90--93.33\%), B+ (86.67--89.99\%), B (83.34--86.66\%), B$-$ (80--83.33\%), C+ (76.67--79.99\%), C (73.34--76.66\%), C$-$ (70--73.33\%), D (60--69.99\%), F ($<$60\%).

\textbf{Functional prototype/solution required}: teams not accomplishing their agreed-upon prototypes or solutions will not receive an A, regardless of other assignment scores. This reflects the reality that engineering is ultimately about producing results, not just following a process.

\section{Course Concerns Chain of Command}
If you have concerns about grading, team dynamics, PA interactions, or course policies, follow this escalation path:
\begin{enumerate}[nosep,leftmargin=1.5em]
\item \textbf{Project Advisor (PA)}: your first point of contact for all project-related concerns.
\item \textbf{Course Lead Departmental Faculty}: if the concern cannot be resolved with your PA.
\item \textbf{Capstone Program Director}: Dr.\ Kristy Csavina (kcsavina@mines.edu). For program-level concerns.
\item \textbf{Dean}: for concerns that cannot be resolved at the program level.
\end{enumerate}

Most concerns are resolved at Step~1. If you skip steps, you will be directed back to the appropriate level.


\section{Engineering Judgment and Professional Responsibility}
Engineering judgment is the ability to make sound decisions when the textbook does not provide a clear answer. It develops through experience, mentorship, and reflection; and capstone is where you begin developing it.

\textbf{When codes and standards do not cover your situation}: engineering standards cannot anticipate every design scenario. When your specific configuration is not addressed by the applicable standard, you must apply the standard's \emph{intent} (protect public safety, ensure reliability, maintain performance) using your engineering knowledge and conservative assumptions. Document your reasoning: ``No standard directly addresses this configuration. We applied the intent of IEC~61010 Section~6.3 by [specific approach], with a safety factor of 2.0 to account for the uncertainty introduced by the non-standard application.''

\textbf{When stakeholder interests conflict}: the client wants the cheapest solution; the end user needs the safest solution; the schedule demands the fastest solution. Engineering judgment means recognizing these conflicts explicitly, evaluating the trade-offs with data, and recommending a course of action that satisfies the most important constraint (safety) while honestly communicating the trade-offs on the others.

\textbf{When you do not know enough}: the ethical response to uncertainty is not to guess confidently; it is to acknowledge the limitation and take appropriate action; research, consult an expert, add safety margin, or design a test to resolve the uncertainty. ``I don't know, but I will find out'' is always preferable to ``I think it's fine'' without evidence.

\section{Generative AI: Detailed Guidance for Capstone}
AI tools (ChatGPT, Claude, Copilot, Gemini, etc.) are powerful engineering aids when used appropriately. The course policy is not anti-AI; it is pro-integrity.

\textbf{Appropriate uses that make you more effective}:
\begin{itemize}[nosep,leftmargin=1.5em]
\item \textbf{Research acceleration}: ``What are the common failure modes of MOSFET H-bridges?'' AI provides a starting point; verify with datasheets and application notes.
\item \textbf{Brainstorming}: ``Generate 10 concepts for a soil moisture sensor mounting mechanism.'' AI augments your ideation; evaluate each concept using your engineering judgment.
\item \textbf{Code generation}: ``Write an ESP32 function to read an I2C temperature sensor.'' AI produces draft code; review, test, and modify to match your specific hardware configuration.
\item \textbf{Proofreading}: ``Check this paragraph for grammar and clarity.'' AI catches errors you miss after staring at the document for hours.
\item \textbf{Learning}: ``Explain the Ziegler-Nichols PID tuning method with an example.'' AI provides a tutorial you can compare against the MEGN~300 reference.
\end{itemize}

\textbf{Inappropriate uses that violate academic integrity}:
\begin{itemize}[nosep,leftmargin=1.5em]
\item Submitting AI-generated text as your own writing in reports, essays, or deliverables without attribution.
\item Using AI to generate ``analysis'' that you do not understand and cannot verify. If you cannot explain the calculation to your PA, you did not do the analysis.
\item Generating citations that do not exist (``hallucinated references''). AI frequently invents plausible-looking citations that refer to papers, books, or websites that do not exist. \textbf{Every citation must be verified by locating the actual source.}
\item Generating ethics essay content that is generic and disconnected from your specific project. The ethics assignment specifically requires engagement with \emph{your} project's unique ethical dimensions; something AI cannot provide without deep project knowledge.
\end{itemize}

\textbf{The attribution rule}: when AI-generated content contributes substantially to a deliverable, disclose it: ``AI assistance was used for [specific purpose, e.g., initial literature search, code template generation, grammar checking] and the output was reviewed and modified by the team.'' Transparency is always acceptable; concealment is always a violation.

\textbf{The verification rule}: AI output is a draft, not a final answer. Every AI-generated claim must be verified against authoritative sources. Every AI-generated calculation must be checked by hand or independent tool. Every AI-generated code snippet must be tested on your hardware. You are responsible for everything in your deliverables, regardless of its origin.



\section{The Broader Impacts Essay}
\begin{keyconceptbox}[Broader Impacts Essay: Quick Reference]
This section provides the complete guide for the SDII Broader Impacts Essay. For the SDI Ethics Assignment (a related but distinct deliverable), see the guidance in Chapter~18 (Professional Communication; Written).
\end{keyconceptbox}

The Broader Impacts Essay (SDII, Sprint~11) is a \textbf{1000--1500 word argumentative essay} on a societal impact connected to your project. It follows the same structure as the Ethics Assignment (Chapter~18) but focuses on broader societal effects rather than narrow ethical dilemmas:

\textbf{Scope}: how does your project's technology, application, or domain affect society beyond the immediate client? Consider: environmental impact, equity and access, labor displacement, privacy and surveillance, public safety, community resilience, economic effects on marginalized populations, and unintended consequences.

\textbf{Structure}: take a position (``Automated agricultural monitoring systems will exacerbate economic inequality among small farmers unless [specific policy intervention]''), defend it with evidence (cite economic data, case studies, policy analyses), address counterarguments honestly, and conclude with a recommendation.

\textbf{Connection to your project}: the essay must be grounded in your specific project's technology and application domain, not a generic discussion of ``technology and society.'' The reader should understand exactly how your project relates to the broader impact you are analyzing.

\textbf{Common mistakes}: (1)~writing an informative essay instead of an argumentative one (no thesis, no position), (2)~discussing impacts so broad they are disconnected from the project (``AI will change everything''), (3)~taking a position so bland that no one would disagree (``access to clean water is important''), (4)~using AI to generate the essay without genuine engagement with the topic.

\section{Academic Integrity in Capstone}
Academic integrity in capstone extends beyond the standard prohibition on plagiarism. Capstone-specific integrity issues:

\textbf{Data fabrication}: presenting test results that were not actually obtained, or modifying data to make results appear more favorable. This is an automatic failing offense and potential expulsion. If a test failed, report the failure honestly and present your analysis of why it failed and how you would address it. Honest failure is infinitely preferable to fabricated success.

\textbf{Contribution misrepresentation}: claiming credit for work you did not perform. Feedback Loop evaluations and PA observations are designed to identify contribution imbalances. If you did not write a section of the report, do not claim authorship. If a teammate completed the analysis, credit them in the document.

\textbf{Unauthorized collaboration}: while capstone is a team project, individual assignments (Ethics Assignment, Broader Impacts Essay, Individual Performance Evaluations) must be your own work. Using a teammate's essay as a template, submitting substantially similar individual assignments, or dividing individual work among team members are all violations.

\textbf{Source attribution}: cite every source. This includes: datasheets, application notes, textbooks, online tutorials, StackOverflow answers, AI-generated content, and prior capstone team reports. When in doubt, cite it. Over-citation is never penalized; under-citation can be an integrity violation.

\textbf{Consequences}: academic integrity violations are adjudicated through the Mines Student Code of Conduct. Consequences range from a zero on the assignment to failure in the course to expulsion from the university, depending on the severity and whether it is a first or repeated offense. Your PA and CF take integrity seriously; do not test the system.


\section{Professional Development Modules}
Each semester includes Professional Development assignments completed on Canvas. These modules build skills that complement the project work:

\textbf{SDI modules} (typical selection): Introduction to Scrum, Creating Team Culture, Intellectual Property, Professional Licensure and Liability, Engineering Ethics (pre-essay preparation), and discipline-specific modules selected by the team.

\textbf{SDII modules} (typical selection): Advanced Scrum, Project Retrospective Techniques, Resume and Interview Preparation, and discipline-specific modules (e.g., CE construction management, EE PCB design best practices, ME FEA validation).

The team's \textbf{Professional Development Plan} identifies which modules the team will complete each semester. The plan is submitted at the start of each semester and reviewed with the PA. Some modules are required for all students; others are selected based on project needs and individual development goals.

\textbf{Individual Professional Development}: CLO~15 requires you to demonstrate functioning effectively as a member of a team. CLO~16 requires applying project management methodologies. CLO~17 requires self-evaluation and self-improvement strategies. These CLOs are assessed through: your Professional Development module completions, your Feedback Loop self-evaluations, your PA's observation of your contributions, and your end-of-semester reflection narrative.

\section{ABET Student Outcome Alignment}
Capstone Design is the primary assessment vehicle for several ABET Student Outcomes~[25]. Your work in capstone directly demonstrates your competence in these areas:

\textbf{SO~1}: an ability to identify, formulate, and solve complex engineering problems by applying principles of engineering, science, and mathematics (demonstrated through: requirements engineering, engineering analysis, design calculations, test methodology).

\textbf{SO~2}: an ability to apply engineering design to produce solutions that meet specified needs with consideration of public health, safety, welfare, global, cultural, social, environmental, and economic factors (demonstrated through: the complete design cycle from problem definition through V\&V, sustainability assessment, ethics analysis).

\textbf{SO~3}: an ability to communicate effectively with a range of audiences (demonstrated through: PDR/CDR/FDR reports and presentations, client communication, Showcase poster and demonstration).

\textbf{SO~5}: an ability to function effectively on a team whose members together provide leadership, create a collaborative and inclusive environment, establish goals, plan tasks, and meet objectives (demonstrated through: team charter, sprint management, Feedback Loop evaluations, PA observation of team dynamics).

\textbf{SO~7}: an ability to acquire and apply new knowledge as needed, using appropriate learning strategies (demonstrated through: self-directed learning of domain-specific technologies, professional development modules, end-of-semester reflection).

Your capstone deliverables are the evidence that you have achieved these outcomes. Strong capstone performance is a direct indicator of your readiness for professional engineering practice.

\section{The Final Course Survey}
At the end of each semester, complete the Final Course Survey on Canvas. This survey provides feedback to the capstone program on: PA effectiveness, course structure, resource availability, and suggestions for improvement. The survey is anonymous and does not affect your grade.

Your feedback matters: the capstone program has evolved significantly based on student survey responses. Specific, constructive feedback (``The CDR timeline was too compressed; starting CDR one week earlier would allow more time for the Engineering Calculation Package'') is far more useful than generic comments (``The course was hard'').





\section{Ethics Beyond the Classroom}
The ethical principles you practice in capstone apply throughout your career:

\textbf{Whistleblowing}: if you discover that your employer is knowingly producing unsafe products, falsifying test data, or violating environmental regulations, you have a professional and legal obligation to report the violation. The path: (1)~raise the concern internally through your chain of command, (2)~if internal channels fail or retaliate, report to the appropriate regulatory agency (OSHA, EPA, SEC, state engineering board). Whistleblower protection laws exist, but the path is difficult. Studying the ethics case studies (Challenger, Flint, Boeing) prepares you to recognize the warning signs before a crisis occurs.

\textbf{Conflicts of interest}: as a PE or engineering employee, you must disclose any financial, personal, or professional interest that could bias your engineering judgment. If your company is evaluating a product made by your spouse's company, you must disclose the relationship and recuse yourself from the evaluation. Undisclosed conflicts of interest can result in loss of PE license, termination, and legal liability.

\textbf{Continuing competence}: the NCEES Rules require that you undertake only assignments for which you are qualified. As technology evolves, the boundary of your competence shifts. An engineer qualified to design analog circuits in 2026 may not be qualified to design AI-based control systems in 2036 without additional training. Continuing education (required for PE license maintenance in most states) ensures that your qualifications keep pace with technology.

\textbf{Pro bono and public service}: many engineering professional societies encourage members to donate engineering services to community organizations, disaster relief, and underserved populations. Engineers Without Borders, Bridges to Prosperity, and similar organizations provide opportunities to apply engineering skills for public benefit. Capstone projects with community clients provide your first experience with this form of professional service.



\begin{examplebox}[Engineering Judgment: The Safety Factor Decision]
Your structural analysis shows that a mounting bracket has a safety factor of 1.8 under the maximum expected load. The requirement specifies SF $\geq$ 2.0. The bracket is machined from 6061-T6 aluminum using conservative material properties from the datasheet.

\textbf{Option A}: redesign the bracket with thicker material to achieve SF = 2.0. Impact: +\$15 material cost, +4 hours machining time, bracket weight increases by 30\%.

\textbf{Option B}: accept SF = 1.8 and document the engineering rationale. The 1.8 value used minimum datasheet properties; typical 6061-T6 yield strength is 10--15\% higher than minimum, which would bring the effective SF above 2.0. Additionally, the ``maximum expected load'' includes a 1.5$\times$ application factor that already accounts for dynamic loading uncertainty.

\textbf{The judgment call}: a strictly by-the-numbers approach says Option~A (the calculated SF is below the requirement). An engineering judgment approach considers: the conservatism already embedded in the analysis (minimum material properties, 1.5$\times$ load factor), the consequence of failure (the bracket supports a sensor module, not a load-bearing structural element), and the schedule impact of redesign (4 hours during a schedule-critical sprint).

\textbf{Professional practice}: document both options. Present the analysis to your PA. If the PA and team agree that the conservatism in the analysis provides adequate margin, Option~B may be acceptable with a documented rationale and client awareness. If the bracket supports a safety-critical function, Option~A is required regardless of schedule impact.

This is engineering judgment: making defensible decisions when the numbers alone do not give a clear answer.
\end{examplebox}


\section{Practice Questions}
\begin{enumerate}[leftmargin=1.5em]
\item Describe the four-step path to PE licensure. At what point in your career should you take the FE exam, and why is it advantageous to do so while still in school?
\item Your client asks you to omit a safety guard from the final design because ``it gets in the way of operation.'' Using the NCEES Rules, explain your professional obligation and how you would handle the conversation.
\item A teammate used ChatGPT to generate a section of the CDR report and submitted it as their own work. The section includes two citations that do not exist. What policies has this violated, and what should happen?
\item Select one of the four case studies (Challenger, Flint, Boeing, Hyatt Regency) and identify the specific process failure that enabled the technical failure. How could a capstone team experience a similar process failure, and what safeguards would prevent it?
\item You receive a grade you believe is unfair. Walk through the chain of command and describe what information you would bring to each conversation.
\end{enumerate}


